\documentclass[11pt,a4paper]{article}
\usepackage[utf8]{inputenc}
\usepackage[T1]{fontenc}
\usepackage[spanish,es-noshorthands,es-nodecimaldot]{babel}
\usepackage{lmodern}
\usepackage{amsmath,amssymb,amsthm}
\usepackage{booktabs}
\usepackage[margin=2.5cm]{geometry}
\usepackage[colorlinks=true,linkcolor=blue,citecolor=blue,urlcolor=blue]{hyperref}

\newtheorem{teorema}{Teorema}
\newtheorem{proposicion}[teorema]{Proposición}
\newtheorem{lema}[teorema]{Lema}
\newtheorem{corolario}[teorema]{Corolario}
\newtheorem{conjetura}[teorema]{Conjetura}
\theoremstyle{remark}
\newtheorem{nota}{Observación}

\newcommand{\G}{\mathcal G}
\newcommand{\Ncal}{\mathcal N}
\newcommand{\dd}{\mathrm d}
\newcommand{\B}{\mathfrak B}
\setlength{\emergencystretch}{3em}

\title{El déficit $d=N$ y el fantasma de la garganta:\\
un mecanismo común para la inestabilidad monopolar de los agujeros de gusano\\
en teorías escalares y vectoriales--tensoriales de segundo orden}
\author{Iván López Martín\\[2pt]
\small Investigador independiente\thanks{Código y datos: \href{https://github.com/iivanlopez/gusanos}{github.com/iivanlopez/gusanos}. Cálculos simbólicos y numéricos realizados con asistencia de Claude (Anthropic).}}
\date{27 de octubre de 2026}

\begin{document}
\maketitle

\begin{abstract}
Estudiamos el sector monopolar ($l=0$) de las gargantas de agujero de gusano estáticas, esféricas y con dos extremos asintóticamente planos en las teorías de segundo orden con simetría de desplazamiento más generales que no han sido descartadas por los teoremas de Horndeski: la teoría cuártica más allá de Horndeski (GLPV), Horndeski quíntico, el escalar lineal en el tiempo $\phi=qt+\psi(\ell)$ y la Proca generalizada cuártica. Mostramos que el monopolo obedece a un mecanismo común con dos ramas. (i)~Si la garganta la sostiene un campo \emph{radial} (un escalar estático o la componente $A_1$ del vector), la función de prueba $\delta r=N$ (el radio areal reescalado con el lapso) cumple todas las ligaduras y tiene energía exactamente
$E[N]=-2\int N^3\G\,[(\ln r/N)']^2\,\dd\ell$, sin funciones libres, donde $\G$ es el coeficiente cinético del sector impar: si este no tiene fantasmas, el monopolo tiene un taquión. (ii)~Si la sostiene un campo \emph{temporal} (el escalar $\phi=qt$ o la componente $A_0$), la energía de las ligaduras puede pagar el déficit, pero el abocinamiento $r''(0)>0$ fija el signo de la energía cinética en la garganta: las dos acciones cinéticas reducidas toman la forma común $T=\int[-(2\G/N)\dot d^2+k((\G\dot d/N)')^2]$ con $k(0)=-\kappa_0/(\G(0)r''(0))$, $\kappa_0>0$, de modo que hay un fantasma monopolar si $\G(0)>0$ y uno tensorial si $\G(0)<0$ (con vector temporal, $2\G$ es el coeficiente $q_T$ de~\cite{DeFeliceProca}). Para el escalar lineal en el tiempo sin término cúbico demostramos además un resultado de existencia más fuerte: $G_4$ cambia de signo en toda garganta. Comprobamos que la rotación lenta, la carga eléctrica, un fluido perfecto y un sólido elástico no cambian el signo del déficit en el sector par, y que las gargantas exactas de Bakopoulos, Charmousis y Kanti tienen el monopolo fuertemente acoplado ($K\equiv0$) con $E[N]<0$. El único caso no cerrado es el galileón cúbico con flujo radial del escalar ($G_3\neq0$, $\psi'\neq0$): damos dos identidades integrales, la forma cerrada del símbolo principal del monopolo, dos lemas nuevos (la garganta es un mínimo de $N G_4$ y la cola debe tener densidad efectiva decreciente) y una criba de unos 3600 fondos en la que ninguno tiene el monopolo localmente bien planteado y sin fantasmas.
\end{abstract}

\tableofcontents

%======================================================================
\section{Introducción}

Un agujero de gusano atravesable estático necesita, en relatividad general (RG), materia que viole la condición de energía nula en la garganta~\cite{MorrisThorne}. El ejemplo canónico, la garganta de Ellis--Bronnikov~\cite{Ellis,Bronnikov} sostenida por un escalar fantasma, es inestable en el sector radial~\cite{GonzalezGuzmanSarbach}. Las teorías escalares--tensoriales de segundo orden ofrecen una alternativa: la violación de la condición nula puede producirla un escalar sano, o un acoplo a la curvatura como en Einstein--escalar--Gauss--Bonnet~\cite{KKK,Canate}. En Horndeski, sin embargo, no hay gargantas estáticas, esféricas y asintóticamente planas libres de fantasmas e inestabilidades de gradiente~\cite{Rubakov,EvseevMelichev}. La demostración se apoya en una estructura que deja de existir más allá de Horndeski~\cite{Franciolini}, y Mironov, Rubakov y Volkova (MRV) construyeron gargantas en la teoría cuártica de Gleyzes--Langlois--Piazza--Vernizzi (GLPV)~\cite{GLPV} que cumplen las condiciones de alta energía en los dos sectores de paridad~\cite{MRV1,MRV2,MRV3}. El sector monopolar se ha tratado como uno más del sector par, y el análisis completo de Mironov y Volkova~\cite{MironovVolkova} supone que no da condiciones adicionales.

Este trabajo resume un programa (las Fases~2f--2m del repositorio~\cite{repo}) que estudia precisamente ese sector, con la motivación de un objeto híbrido: una estrella con una garganta en su centro. El resultado es negativo y, sobre todo, \emph{estructural}: el monopolo de todas las gargantas de esas teorías falla por un mismo mecanismo, que solo depende de la geometría de dos extremos y del coeficiente cinético del sector impar o tensorial. Su forma depende de qué campo sostiene la garganta:
\begin{itemize}
\item \textbf{campo radial} (un escalar estático, o la componente radial $A_1$ de un vector): una función de prueba universal, $\delta r=N$, tiene energía negativa sin funciones libres (\S\ref{sec:radial});
\item \textbf{campo temporal} (un escalar lineal en el tiempo, o la componente $A_0$): la energía de las ligaduras puede compensar ese déficit, pero el abocinamiento de la garganta invierte el signo de la energía cinética (\S\ref{sec:temporal}).
\end{itemize}
Las salidas naturales (rotación, carga, materia, el término quíntico, gargantas exactas por transformación disforme) no cambian la conclusión (\S\ref{sec:salidas}). Queda un caso sin teorema, el galileón cúbico con flujo radial, para el que damos lemas exactos y evidencia numérica (\S\ref{sec:cubico}).

\paragraph{Convenciones.} Signatura $(-,+,+,+)$, $X=-\tfrac12(\partial\phi)^2$, $G=c=1$ salvo factores $16\pi$ absorbidos en las funciones de la teoría. El fondo es
\begin{equation}
\dd s^2=-N(\ell)^2\dd t^2+\dd\ell^2+r(\ell)^2\dd\Omega^2,\qquad N=e^{\Phi},
\label{eq:fondo}
\end{equation}
con $\ell\in\mathbb R$ la longitud propia, una garganta en $\ell=0$ ($r'(0)=0$, $r''(0)>0$) y $r\to\infty$, $N\to1$ en los dos extremos. Las perturbaciones $l=0$ son $\delta g_{tt}=-2N^2a$, $\delta g_{t\ell}=\beta$, $\delta g_{\ell\ell}=2c$, $\delta r=d$.

%======================================================================
\section{Teorías y estructura del monopolo}\label{sec:marco}

\paragraph{Escalar.} La GLPV cuártica~\cite{GLPV,Kobayashi},
\begin{equation}
\mathcal L=G_2-G_3\Box\phi+G_4R+G_{4X}\big[(\Box\phi)^2-\phi_{\mu\nu}^2\big]+F_4\,\epsilon^{\mu\nu\rho}{}_{\sigma}\epsilon^{\mu'\nu'\rho'\sigma}\phi_\mu\phi_{\mu'}\phi_{\nu\nu'}\phi_{\rho\rho'},
\label{eq:GLPV}
\end{equation}
con $G_i(\phi,X)$ y $F_4(\phi,X)$, y su extensión quíntica de Horndeski con $G_5$. Definimos
\begin{equation}
\G=G_4-2XG_{4X}-4X^2F_4\qquad\big(\text{quíntico: }\G=G_4-2XG_{4X}+X(\Phi'\phi'G_{5X}+G_{5\phi})\big).
\label{eq:Gcal}
\end{equation}
$\G$ es el coeficiente cinético del sector impar~\cite{KMS1} (con $\G>0$, el sector impar no tiene fantasmas) y también, como veremos, el de las ondas tensoriales en el corte $\phi=$~cte.

\paragraph{Vector.} La Proca generalizada cuártica~\cite{Heisenberg},
\begin{equation}
\mathcal L=G_2(X,F)+G_3(X)\nabla_\mu A^\mu+G_4(X)R+G_{4X}\big[(\nabla_\mu A^\mu)^2-\nabla_\rho A_\sigma\nabla^\sigma A^\rho\big],
\end{equation}
con $A_\mu=(A_0,A_1,0,0)$, $X=-\tfrac12A_\mu A^\mu$ y $F=-\tfrac14F_{\mu\nu}F^{\mu\nu}$. Aquí $\G=G_4-2XG_{4X}$.

\paragraph{Maquinaria.} Para cada teoría derivamos simbólicamente la acción reducida esférica y su desarrollo cuadrático $l=0$ con funciones de la teoría generales, (para el sector par de Horndeski, véase~\cite{KMS2}) y lo comprobamos con identidades de Noether exactas en aritmética racional, con modos gauge exactos sobre chorros en la capa de masas y con soluciones conocidas: Schwarzschild, Reissner--Nordström, Kerr lineal, Hartle--Thorne, el agujero negro de Lü--Pang~\cite{LuPang} a 40 cifras y el agujero negro ``stealth'' de Babichev--Charmousis~\cite{BabichevCharmousis} a 50 cifras.

\paragraph{Estructura del monopolo con escalar estático.} En el gauge unitario $\delta\phi=0$ (Fases~2f y 2j), exactamente,
\begin{equation}
L_2=-F\,W+L_{\rm rest}(c,d),\qquad F=N^2a'+\dot\beta,\qquad W=\gamma c+\alpha_1d'+\alpha_0d,\qquad \alpha_0=-\Phi'\alpha_1 .
\label{eq:estructura}
\end{equation}
El lapso y el desplazamiento entran solo a través del multiplicador $F$, que impone la ligadura de momento $W=0$; $c$ es auxiliar; tras eliminarla queda un único campo $d$ con $S_{\rm red}=\int\dd t\,\dd\ell\,[K\dot d^2-Pd'^2-Ud^2]$. La identidad $\alpha_0=-\Phi'\alpha_1$ es exacta en la cuártica y en la quíntica.

%======================================================================
\section{Campo radial: el déficit $d=N$}\label{sec:radial}

\begin{teorema}[déficit del monopolo; Fases 2g, 2j, 2m]\label{teo:dN}
Sea un fondo~\eqref{eq:fondo} con escalar estático $\phi(\ell)$ de la GLPV cuártica~\eqref{eq:GLPV} o de Horndeski quíntico, o con vector puramente radial ($A_0=0$) de la Proca generalizada cuártica, con funciones de la teoría arbitrarias. La configuración
\begin{equation}
d=N f,\qquad c=0,\qquad \delta\phi=0\ (\text{resp. }\delta A=0)
\end{equation}
cumple todas las ligaduras, y su energía estática es, punto a punto y salvo derivadas totales,
\begin{equation}
E[Nf]=-2\int N^3\,\G\,\big(f'-\sigma f\big)^2\,\dd\ell,\qquad \sigma\equiv\Big(\ln\frac rN\Big)' .
\label{eq:EN}
\end{equation}
En particular, $E[N]=-2\int N^3\G\sigma^2\,\dd\ell<0$ si $\G>0$. Si además el monopolo está bien planteado ($K,P>0$), existe un modo con $\omega_0^2\le E[N]/\int KN^2\,\dd\ell<0$.
\end{teorema}

\begin{proof}[Esquema]
Como $\alpha_0=-\Phi'\alpha_1$, la ligadura es $W=\gamma c+\alpha_1N(d/N)'$, que se anula con $c=0$ y $d\propto N$ (para $f$ constante; para $f$ general se corrige con $c$). La energía de la prueba es la parte de $L_{\rm rest}$ que solo contiene $d$, cuyos coeficientes, en la capa de masas, solo dependen de $\G$; la identidad~\eqref{eq:EN} se comprueba simbólicamente fuera de la capa de masas. En la quíntica, la reducción a la forma canónica (única) da el coeficiente de Kobayashi--Motohashi--Suyama~\eqref{eq:Gcal}. Para el vector radial, el fondo es el de Horndeski con $\phi'=A_1$, y en el gauge $\delta A_1=0$ la componente $\delta A_0$ solo se acopla a velocidades. Por último, $\sigma\not\equiv0$ porque $r/N\to\infty$ en los extremos, y $d=N\to1$ es normalizable porque $K$ es integrable.
\end{proof}

\begin{nota}
El cociente de Rayleigh no contiene $G_2$, $G_3$ ni la normalización del escalar: el teorema vale en campo fuerte y débil, con cualquier masa. En campo débil ($|\Phi'|r\ll1$) el taquión es \emph{rápido}: $\omega_0^2\lesssim-O(1)/(Mr_0)$ (Fase~2f). En una garganta con la masa del Sol y $r_0\sim1$~km, el tiempo de crecimiento es de milisegundos.
\end{nota}

\paragraph{Invariancia de marco.} Con escalar estático, una transformación disforme $\bar g_{\mu\nu}=C(X)g_{\mu\nu}+D(X)\phi_\mu\phi_\nu$ deja invariante la prueba ($\delta X=0$ y el lapso no cambia), de modo que su energía es la misma en los dos marcos (Fase~2j). Esto lleva el teorema a las gargantas exactas de Bakopoulos--Charmousis--Kanti (BCK)~\cite{BCK}, construidas por transformación disforme a partir del agujero negro de Lü--Pang (\S\ref{sec:BCK}).

%======================================================================
\section{Campo temporal: el fantasma de la garganta}\label{sec:temporal}

Si el escalar depende linealmente del tiempo, $\phi=qt+\psi(\ell)$, la métrica puede ser estática en teorías con simetría de desplazamiento~\cite{BabichevCharmousis}. El gauge unitario fija entonces el \emph{tiempo}, y la estructura~\eqref{eq:estructura} se pierde: el lapso y el desplazamiento del corte $\phi=$~cte pasan a ser auxiliares algebraicos, con términos cuadráticos propios, y la prueba $d=N$ deja de cumplir las ligaduras. Lo mismo ocurre con la componente temporal $A_0$ de un vector, que en el lenguaje de Stückelberg es un $\dot\pi$.

\subsection{El escalar lineal en el tiempo}

Con $\tau=t+\psi/q$ el escalar es $\phi=q\tau$, y en las variables ADM del corte (lapso $\Ncal$, desplazamiento $v$, longitud propia $\rho$) la lagrangiana es la de la teoría efectiva~\cite{GLPV}
\begin{equation}
L=A_2+A_3K+A_4\big(K^2-K_{ij}K^{ij}\big)+B_4R^{(3)},\qquad A_4=-G_4+2XG_{4X}+4X^2F_4=-\G,\quad B_4=G_4 ,
\end{equation}
con $A_2=G_2$ y $A_3=yG_3-\int_0^yG_3\,\dd y$, $y=\sqrt{2X}$. La velocidad tensorial en el corte es $c_T^2=B_4/\G$. Para configuraciones estacionarias obtenemos en forma cerrada (Fase~2k; la diferencia con la lagrangiana covariante tiene operador de Euler--Lagrange idénticamente nulo)
\begin{equation}
\mathcal L_{\rm est}=\Ncal r^2A_2-A_3(r^2v)'+\frac{2A_4}{\Ncal}\,r'(v^2r)'+\Ncal B_4\big(2-4rr''-2r'^2\big),
\end{equation}
con la ligadura de momento $rA_3'=2v\,(Qr')'$, $Q\equiv2A_4/\Ncal$. Sin término cúbico ($A_3'\equiv0$ en el fondo), una garganta obliga a $v\equiv0$ ($\psi'=0$) o a $A_4\equiv0$.

\begin{teorema}[no-go de existencia; Fase 2k]\label{teo:beta}
Sea una garganta estática con $\phi=qt+\psi$ de la GLPV cuártica con simetría de desplazamiento, con $v\equiv0$ o $A_3'\equiv0$, dos extremos asintóticamente planos ($r'/r^2\to0$ y $\beta'/(r\beta)\to0$) y $G_4\to G_{4\infty}\neq0$. Entonces $\beta\equiv\Ncal G_4$ cambia de signo; como $\Ncal>0$, $G_4$ cambia de signo en el fondo.
\end{teorema}

\begin{proof}
Eliminando $A_2$ entre la primera integral y la ecuación radial, $\beta$ cumple, sin ninguna función libre,
\begin{equation}
\Big(\frac{\beta'}{r}\Big)'+\frac{K\beta}{r^3}=0,\qquad K=1+rr''-r'^2 .
\label{eq:beta}
\end{equation}
Si $\beta$ no se anula, dividiendo por $\beta$ e integrando por partes, y usando $K/r^3=(r'/r^2)'+(1+r'^2)/r^3$,
\begin{equation}
0\le\int_{\mathbb R}\frac{\beta'^2}{r\beta^2}\,\dd\rho=\Big[\frac{\beta'}{r\beta}\Big]-\int\frac{K}{r^3}\,\dd\rho=-\Big[\frac{r'}{r^2}\Big]-\int\frac{1+r'^2}{r^3}\,\dd\rho<0. \qedhere
\end{equation}
\end{proof}

No hace falta simetría de reflexión ni unicidad de la garganta; lo esencial es la topología de dos extremos (con un centro regular, $r(0)=0$ y la integral no existe). Donde $B_4=G_4$ se anula, $c_T^2=0$ (acoplo fuerte); a un lado del cero hay una inestabilidad de gradiente tensorial. En RG ($\beta=N$), \eqref{eq:beta} es la condición de presión isótropa.

\paragraph{El monopolo con $\phi=qt$.} La acción cuadrática dependiente del tiempo (exacta salvo derivadas totales) tiene el lapso solo en la parte estática y el desplazamiento solo en la cinética. Eliminándolos,
\begin{equation}
T[\dot d]=\int\Big[Q\dot d^2+\frac{r\,[(Q\dot d)']^2}{(Qr')'}\Big]\dd\rho,\qquad
V[d]=\int\Big[\frac{S^2}{2\kappa}-2\beta r^2\Big[\Big(\frac dr\Big)'\Big]^2\Big]\dd\rho,
\label{eq:Tqt}
\end{equation}
sin términos giroscópicos. La familia $\delta X=0$, análoga de la prueba $d=N$, tiene energía $-2\int\beta r^2[(d/r)']^2$: la misma forma que~\eqref{eq:EN}, con $\G\to G_4$, y el Teorema~\ref{teo:beta} impide precisamente que $\beta>0$.

\subsection{El vector temporal}

Con $A_1=0$ y $G_3=0$ (Fase~2m), la estructura exacta es $L_2=\kappa F^2-FW+L_{\rm rest}$ con $\kappa=U^2r^2(G_{2F}+2FG_{2FF})/(2N^5)$ ($U$ es $A_0$ en el fondo). La prueba $d=N$ no cumple la ligadura, y eliminando $F$ y la perturbación algebraica de $A_0$ aparece la energía del campo eléctrico perturbado,
\begin{equation}
E[N]=\int\Big[-2N^3G_4\sigma^2+\frac{2N}{g}\big(2G_{4X}N(U/N)'-G_{2F}U'\big)^2+\frac{J_N^2}{4\tilde m}\Big]\dd\ell,\qquad g=G_{2F}+2FG_{2FF},
\end{equation}
que en las seis gargantas regulares que construimos supera al déficit por un factor de 1.2 a 4.6. Pero la garganta impone una condición local exacta,
\begin{equation}
r''(0)=-\frac{X\Theta}{2r_0\G},\qquad \tilde c(0)=-\frac{N_0\Theta}{2}=\frac{N_0r_0\G}{X}\,r''(0),\qquad \Theta\equiv G_{2X}r_0^2+2G_{4X},
\end{equation}
donde $\tilde c$ es el coeficiente de la componente radial auxiliar $\tilde w=\delta A_1+U\beta/N^2$. La reducción completa en el gauge $c=0$ da
\begin{equation}
S_{\rm red}=\int\Big[-\frac{2\G}{N}\dot d^2+k_2\Big(\Big(\frac{\G\dot d}{N}\Big)'\Big)^2-V(d)\Big],\qquad k_2=-\frac{4r^2N^4}{U^2\tilde c}.
\end{equation}

\subsection{La regla de signos común}

\begin{teorema}[fantasma de la garganta]\label{teo:fantasma}
En los dos casos temporales sin flujo ($\phi=qt$ con $\psi'=0$; vector con $A_1=0$, $G_3=0$) la energía cinética del monopolo es
\begin{equation}
T=\int\Big[-\frac{2\G}{N}\,\dot d^2+k\,\Big(\Big(\frac{\G\dot d}{N}\Big)'\Big)^2\Big]\dd\ell,\qquad k(0)=-\frac{\kappa_0}{\G(0)\,r''(0)},\quad \kappa_0>0,
\label{eq:Tcomun}
\end{equation}
con $k=4r/(Qr')'$, $\kappa_0=2N_0r_0$ para el escalar y $k=k_2$, $\kappa_0=4N_0^3Xr_0/U_0^2$ para el vector. Por tanto, cerca de una garganta ($r''(0)>0$) los dos términos cinéticos tienen el signo de $-\G(0)$: si las ondas tensoriales del corte no son fantasmas ($\G(0)>0$), toda perturbación $l=0$ con soporte cerca de la garganta tiene energía cinética negativa; si $\G(0)<0$, el fantasma está en el sector tensorial.
\end{teorema}

\begin{proof}
Para el escalar, $Q=2A_4/N=-2\G/N$ transforma~\eqref{eq:Tqt} en~\eqref{eq:Tcomun} con $k=4r/(Qr')'$, y en la garganta $(Qr')'=Qr''$. Para el vector, basta sustituir $\tilde c(0)$ en $k_2$. Las dos sustituciones se comprueban simbólicamente (\texttt{fase2n/regla\_comun.py}).
\end{proof}

Con $G_4=1$, $(Qr')'\propto\rho+p$ es la combinación de la condición nula, y~\eqref{eq:Tcomun} es la energía cinética de un fluido: el teorema es la versión general de ``fluido fantasma $\Rightarrow$ fantasma''. La precisión del símbolo principal en la garganta (Fase~2k \S12) es que, con $\G(0)>0$, o las velocidades radiales son imaginarias (problema mal planteado) o el sector entero tiene el signo cambiado frente a los gravitones. Con conos causales sin inclinar en la garganta, el criterio de corte de Babichev et al.~\cite{BCEL1,BCEL2} no lo evita.

%======================================================================
\section{Las salidas que no funcionan}\label{sec:salidas}

\subsection{Carga y rotación lenta}

Un campo eléctrico radial (``carga sin carga'') no añade grados de libertad en $l=0$ y suma una presión favorable, $E[N]=\int[-2N^3\G\sigma^2+8q_e^2N^3/r^4]$. Pero en una garganta simétrica el coeficiente de la ecuación de $c$ deja de anularse, $e_0=4q_e^2N_0/r_0^3$, lo que fuerza $\delta r(0)=0$ y produce una caída al centro, $U\simeq-C/\ell^2$ con $C/P(0)=7.5$--$19$ (Fase~2g).

En rotación lenta (Fase~2l), la reducción de Kaluza--Klein de la GLPV cuártica da exactamente $L_\omega=\tfrac12\mathcal J\omega'^2$ con $\mathcal J=\tfrac23C^2\G/\sqrt{AB+b^2}$, el mismo $\G$. A momento angular fijo, el sector $l=0$ a $O(\Omega^2)$ es un problema de Routh con $L_{\rm eff}=L_{l=0}-\pi^2/(2\mathcal J)$, y
\begin{equation}
E_J[Nf]=\int\Big[-2N^3\G(f'-\sigma f)^2+\frac{9\pi^2N^3}{\G r^6}f^2\Big]\dd\ell .
\end{equation}
Con $v_0$ la velocidad del observador de momento angular nulo en la garganta, $E_J[N]/|E_0|=-1+\kappa v_0^2$: en Ellis--Bronnikov $\kappa=24/\pi^2$ y el signo solo cambia en $v_0=\pi/\sqrt{24}\approx0.64$; en 27 gargantas reconstruidas, $v_c=0.59$--$1.74$. En el régimen controlado ($v_0\le0.3$) la rotación reduce $|E[N]|$ como mucho un 26\,\%. Además reintroduce la obstrucción de la carga en la garganta, salvo en la subclase $(\G-\phi'^2\G_X)(0)=0$.

\subsection{Materia: fluido y sólido}

Con un fluido perfecto (Fases~2g y 2h), la prueba con el fluido quieto da $E=-2\int N^3[\G\sigma^2-16\pi\Gamma_1p]$, y en el Sol $16\pi\Gamma_1pr_0^2\sim10^{-20}$--$10^{-8}$: el híbrido solar no se salva. Los híbridos densos que cumplían las condiciones necesarias con la superficie de la estrella fija fallan con la superficie libre: el desplazamiento \emph{sin compresión} $y$ tiene
\begin{equation}
E_{Q0}=-2\int_0^{\ell_s}N^3\G(y'-\sigma y)^2\,\dd\ell+E_{\rm out}\,y(\ell_s)^2+E_s\,y'(\ell_s)^2 ,
\end{equation}
con los dos primeros términos negativos. Es un teorema para superficies blandas ($P\propto(Y-Y_s)^m$, $1<m<2$); con salto de densidad, el mejor margen en unos 3200 fondos optimizados es $-0.22$. La velocidad del modo rápido del fluido diverge además en la garganta.

Un sólido elástico (Fase~2i; formulación de campo equivalente a la de~\cite{KarloviniSamuelsson}) estabiliza el sector impar del monopolo (de $-0.34$ a $+0.06$ con $c_T^2=0.2$), pero no el par: en toda región donde el fondo cumple la relación de RG $\Phi'r'-r''=r(\tilde\rho+\tilde p_r)/(4\G)$, la traslación uniforme $d=r'C$, $\xi=-C$ no deforma el material y cumple la ligadura, con energía
$E_m=C^2N[\tilde p_t(r'^2+rr'')+\tilde\rho\,r^2N''/(2N)]$ \emph{sin ningún módulo elástico}. El mejor margen en unos 2900 fondos es $-0.209$.

\subsection{Gargantas BCK}\label{sec:BCK}

En el marco de la semilla (el exterior del agujero negro de Lü--Pang, duplicado), $E[N]<0$ en los nueve juegos de parámetros probados ($-0.17$ a $-0.53$ en unidades de $M$), con $\G\ge1$. Además, el coeficiente cinético del monopolo se anula idénticamente, $K\equiv0$, sobre esa solución (no en fondos quínticos genéricos), y $P<0$ en alguna región: el monopolo de BCK está fuertemente acoplado y el análisis lineal no se aplica.

%======================================================================
\section{El caso abierto: galileón cúbico con flujo}\label{sec:cubico}

Con $G_3\neq0$ y $\psi'\neq0$ la ligadura de momento permite $v\neq0$ con $A_4\neq0$, y~\eqref{eq:beta} gana una fuente:
\begin{equation}
\Big(\frac{\beta'}r\Big)'+\frac{K\beta}{r^3}=-\frac{(Qr^4u')'}{4r^3},\qquad u=\frac{v^2}{r^2}.
\label{eq:betafuente}
\end{equation}
Una configuración simétrica tiene $v$ impar. Las condiciones de salud del fondo son $G_4>0$, $A_4<0$ (es decir $\G>0$, $Q<0$), $\Ncal>|v|$ y $G_4\to1$, $A_4\to-1$ en los extremos.

\paragraph{No hay no-go algebraico.} Dos identidades exactas de~\eqref{eq:betafuente}:
\begin{align}
&\text{(I1)}\quad \int\frac{K\beta}{r^3}\,\dd\rho=\frac34\int\Big(\frac vr\Big)^2(Qr')'\,\dd\rho,\\
&\text{(I2)}\quad \int\frac1r\Big(\frac{\beta'}\beta+\frac{Qr^2u'}{8\beta}\Big)^2+\int\frac{1+r'^2}{r^3}=-\frac34\int\frac{Qu'r'}\beta+\frac1{64}\int\frac{Q^2r^3u'^2}{\beta^2}\qquad(\beta>0).
\end{align}
El último término de (I2) es positivo: un flujo bastante fuerte paga la geometría de dos extremos, y hay fondos explícitos con $G_4>0$ y $A_4<0$ en todo el eje.

\paragraph{Símbolo principal.} La lagrangiana ADM dependiente del tiempo con desplazamiento es exacta frente a la covariante, y el símbolo principal del monopolo tiene forma cerrada. En un punto cualquiera, el coeficiente cinético en el corte $\tau$ es
\begin{equation}
s_2=-\frac{2\B}{r\,v^2r''^2},\qquad \B=-2\beta'(rK)'-Q\,w\,r''(r^4w')'+Q'\,r^4w'(r'w'-r''w),\qquad w=\frac vr ,
\label{eq:B}
\end{equation}
independiente del lapso. En la garganta solo depende de datos de Taylor del fondo:
\begin{equation}
\hat\Lambda=-\frac{\Ncal_0}{4R_0^3}\Big[Q_0^2R_0^4V_1^4+2Q_0R_0^3R_2V_1^2\beta_0-4\beta_0^2(1+R_0R_2)^2\Big],
\end{equation}
con $R_k=r^{(k)}(0)$, $V_k=v^{(k)}(0)$; el gradiente es sano si $\hat\Lambda/A_4>0$.

\begin{lema}[la garganta es un mínimo de $\beta$; Fase 2n]\label{lem:minimo}
Si $\beta_0>0$, $Q_0<0$ y el gradiente del monopolo es sano en la garganta ($\hat\Lambda<0$), entonces $\beta''(0)>R_2\beta_0/(2R_0)>0$.
\end{lema}
\begin{proof}
De~\eqref{eq:betafuente} en $\rho=0$, $\beta''(0)=-(\beta_0/R_0^2)\,[1+R_0R_2+x/2]$ con $x=Q_0R_0^2V_1^2/\beta_0<0$. El corchete de $\hat\Lambda$ es $\beta_0^2[x^2+2R_0R_2x-4(1+R_0R_2)^2]$, positivo con $x<0$ solo si $-x>R_0R_2+\sqrt{R_0^2R_2^2+4(1+R_0R_2)^2}>R_0R_2+2(1+R_0R_2)$. Sustituyendo, $\beta''(0)R_0^2/\beta_0>R_0R_2/2$.
\end{proof}

\begin{lema}[la cola; Fase 2n]\label{lem:cola}
(a) Para cualquier $r(\rho)$, $r^3\big[(r''/r)'+Kr'/r^3\big]=(rK)'$, es decir, $r'$ resuelve~\eqref{eq:beta} si y solo si $rK$ es constante. (b) Con la masa de Misner--Sharp del corte, $r'^2=1-2m(r)/r$, se tiene $K=3m/r-m_r$ y
\begin{equation}
(rK)'=-4\pi\,r'\,r^3\,\frac{\dd\varepsilon}{\dd r},\qquad \varepsilon\equiv\frac{m_r}{4\pi r^2}.
\end{equation}
(c) Donde el flujo se desvanece ($v,v',v''\to0$), $\B\to-2\beta'(rK)'$, de modo que un coeficiente cinético sano exige $\beta'\,\dd\varepsilon/\dd r<0$. Si además el flujo es despreciable más allá (donde $\beta'/r=\int_\rho^\infty K\beta/r^3>0$), la densidad efectiva $\varepsilon$ debe ser decreciente.
\end{lema}

El Lema~\ref{lem:cola} excluye, con cualesquiera $A_4$ y lapso, todas las colas de tipo Ellis ($\varepsilon<0$ creciente hacia cero, $(rK)'=-2r'/r^2$). Obsérvese que la condición cinética del régimen sin flujo, $(Qr')'>0$, sí puede cumplirse en una cola de Ellis con un perfil de cola de $A_4$; es el régimen de número de onda alto el que la excluye. Por otra parte, la dependencia afín de $\B$ en $Q'$ impide cerrar el problema con una identidad lineal $\int\mu\B$ de peso fijo: tras integrar $Q'$ por partes, el integrando $Q\,C$ contiene el término $-2\mu r^4r'w'w''$, de signo indefinido.

\paragraph{Evidencia numérica.} Reconstruimos fondos eligiendo $(r,v,Q,\Ncal)$ y resolviendo~\eqref{eq:betafuente} para la solución par acotada, y evaluamos el símbolo principal en 420--750 puntos por fondo:
\begin{center}
\small\begin{tabular}{p{6.2cm}ccc}
\toprule
familia & fondos & $\beta>0$ & localmente sanos \\
\midrule
Ellis, uno o dos flujos, $A_4$ con pico (2k) & 872 & 56 & 0 \\
cola de energía positiva, $\varepsilon>0$ en todo el eje (2k) & 400 (+ optim.) & 0 & 0 \\
núcleo de Ellis y cola de energía positiva (2n) & 600 (+ 2160 optim.) & 34 & 0 \\
\bottomrule
\end{tabular}
\end{center}
Los fondos fallan por una capa no hiperbólica junto a la garganta (donde $\B$ cambia de signo) o por una cola donde el coeficiente cinético es negativo. La clase con núcleo de Ellis y cola de densidad decreciente, la que deja el Lema~\ref{lem:cola}, admite $\beta>0$ con cola sana (19 de 600 fondos), en contra de lo que sugerían las familias con $\varepsilon>0$ en la garganta. La evolución diferencial en esa clase encuentra un fondo hiperbólico en todo el eje, con $\beta>0$, garganta y cola sanas, pero con una capa fantasma en $0.56\lesssim\rho/r_0\lesssim1.7$: allí $s_2<0$ y $s_0>0$ a la vez. Las dos velocidades características tienen signos opuestos (su producto es $s_0/s_2<0$), así que el corte $\tau$ es espacial para el monopolo y la energía negativa es un fantasma también con el criterio de~\cite{BCEL2}. Es la región donde el pico de $|A_4|$ que sostiene la garganta se deshace ($Q'>0$) y el término afín en $Q'$ de~\eqref{eq:B} tiene el signo malo.

\begin{conjetura}\label{conj:cubico}
Toda garganta estática, simétrica y asintóticamente plana del galileón cúbico con flujo en la GLPV cuártica con simetría de desplazamiento tiene un monopolo mal planteado o con fantasma.
\end{conjetura}

%======================================================================
\section{Discusión}\label{sec:discusion}

\paragraph{El mecanismo.} En los dos casos la causa es la topología de dos extremos. Con un campo radial, la garganta paga con $\sigma=(\ln r/N)'\neq0$: el cociente $r/N$ tiene que crecer hacia los dos lados, y la prueba $d=N$ lo convierte en energía negativa ponderada por el mismo $\G$ que hace sano el sector impar. Con un campo temporal, el potencial de Sturm--Liouville integrado de~\eqref{eq:beta}, $\int(1+r'^2)/r^3$, es positivo, y el abocinamiento $r''(0)>0$ liga el signo de la energía cinética al de $-\G(0)$. Ningún término de la teoría puede compensar el primero (no contiene funciones libres); el segundo solo se compensa invirtiendo el sector tensorial.

\paragraph{Relación con la literatura.} El teorema~\ref{teo:dN} completa por el sector monopolar las condiciones de MRV y de Mironov--Volkova~\cite{MRV3,MironovVolkova}: sus gargantas de alta energía sana tienen un taquión en $l=0$. Para el escalar lineal en el tiempo, la inestabilidad de los agujeros negros con pelo~\cite{OgawaKobayashiSuyama} y el acoplo fuerte de las soluciones stealth~\cite{deRhamZhang} tienen aquí su análogo en el cambio de signo de $G_4$ (Teorema~\ref{teo:beta}). Para la Proca generalizada, el resultado es coherente con las inestabilidades conocidas de soluciones con pelo vectorial~\cite{HKMT,KaseOdd,DeFeliceCubic,GarciaSaenz} y con la observación de que las gargantas vectoriales mínimas necesitan un fantasma~\cite{Ganiyeva}.

\paragraph{Lo que no cubre.} (i) El galileón cúbico con flujo (Conjetura~\ref{conj:cubico}) y el caso vectorial mixto $A_0,A_1\neq0$, que es el mismo hueco y además rompe la paridad temporal (hay términos giroscópicos). (ii) La reducción ADM quíntica con $\phi=qt+\psi$. (iii) DHOST degenerado general fuera de las clases conectadas por transformaciones disformes con escalar estático. (iv) La rotación relativista ($v_0\gtrsim0.6$), fuera de la teoría de perturbaciones. (v) Materia acoplada directamente al escalar, que podría ``sentirlo'' en todo su volumen e impedir la traslación casi gauge.

\paragraph{Consecuencia para el híbrido estrella--garganta.} La garganta solar es de campo extremadamente débil ($M/r_0\sim10^{-6}$): por el Teorema~\ref{teo:dN} tiene un taquión con tiempo de crecimiento $\sim\sqrt{Mr_0}$, de milisegundos, y la materia de la estrella no lo detiene.

%======================================================================
\section{Conclusiones}

\begin{enumerate}
\item Con un campo radial, la prueba $d=N$ da $E[N]=-2\int N^3\G\sigma^2<0$ sin funciones libres en la GLPV cuártica, en Horndeski quíntico y en la Proca generalizada radial: taquión si el sector impar es sano.
\item Con un campo temporal sin flujo, la energía cinética del monopolo toma la forma común~\eqref{eq:Tcomun} y la garganta fija su signo: fantasma monopolar si $\G(0)>0$, tensorial si $\G(0)<0$. Para el escalar, además, $G_4$ cambia de signo.
\item La rotación lenta, la carga, el fluido y el sólido no cambian el signo del sector par; las gargantas BCK están fuertemente acopladas.
\item Queda abierto el galileón cúbico con flujo, con lemas exactos que restringen la garganta y la cola, y unos 3600 fondos sin ningún monopolo localmente sano; el mejor es hiperbólico, pero con una capa fantasma entre la garganta y el máximo del flujo.
\end{enumerate}

\appendix
\section{Reproducibilidad}
Todos los cálculos están en~\cite{repo}: la maquinaria simbólica y las comprobaciones de cada fase en \texttt{fase2f/}--\texttt{fase2m/}, y los lemas de la \S\ref{sec:cubico} y la regla de signos común en \texttt{fase2n/} (\texttt{lemas.py}, \texttt{regla\_comun.py}, \texttt{criba\_mixta.py}, \texttt{optimiza\_mixta.py}). Las referencias de arXiv se verificaron con su API (\texttt{fase2n/arxiv\_bib.py}).

\begin{thebibliography}{99}

\bibitem{MorrisThorne} M.~S. Morris y K.~S. Thorne, \emph{Wormholes in spacetime and their use for interstellar travel: A tool for teaching general relativity}, Am. J. Phys. \textbf{56} (1988) 395.

\bibitem{Ellis} H.~G. Ellis, \emph{Ether flow through a drainhole: A particle model in general relativity}, J. Math. Phys. \textbf{14} (1973) 104.

\bibitem{Bronnikov} K.~A. Bronnikov, \emph{Scalar-tensor theory and scalar charge}, Acta Phys. Pol. B \textbf{4} (1973) 251.

\bibitem{GonzalezGuzmanSarbach} J.~A. González, F.~S. Guzmán y O. Sarbach, \emph{Instability of wormholes supported by a ghost scalar field. I. Linear stability analysis}, Class. Quant. Grav. \textbf{26} (2009) 015010, \href{https://arxiv.org/abs/0806.0608}{arXiv:0806.0608}.

\bibitem{Rubakov} V.~A. Rubakov, \emph{Can Galileons support Lorentzian wormholes?}, \href{https://arxiv.org/abs/1509.08808}{arXiv:1509.08808}.

\bibitem{EvseevMelichev} O.~A. Evseev y O.~I. Melichev, \emph{No static spherically symmetric wormholes in Horndeski theory}, Phys. Rev. D \textbf{97} (2018) 124040, \href{https://arxiv.org/abs/1711.04152}{arXiv:1711.04152}.

\bibitem{Franciolini} G. Franciolini, L. Hui, R. Penco, L. Santoni y E. Trincherini, \emph{Stable wormholes in scalar-tensor theories}, JHEP \textbf{01} (2019) 221, \href{https://arxiv.org/abs/1811.05481}{arXiv:1811.05481}.

\bibitem{GLPV} J. Gleyzes, D. Langlois, F. Piazza y F. Vernizzi, \emph{Healthy theories beyond Horndeski}, Phys. Rev. Lett. \textbf{114} (2015) 211101, \href{https://arxiv.org/abs/1404.6495}{arXiv:1404.6495}.

\bibitem{MRV1} S. Mironov, V. Rubakov y V. Volkova, \emph{Towards wormhole beyond Horndeski}, EPJ Web Conf. \textbf{191} (2018) 07014, \href{https://arxiv.org/abs/1811.05832}{arXiv:1811.05832}.

\bibitem{MRV2} S. Mironov, V. Rubakov y V. Volkova, \emph{More about stable wormholes in beyond Horndeski theory}, \href{https://arxiv.org/abs/1812.07022}{arXiv:1812.07022}.

\bibitem{MRV3} S. Mironov, V. Rubakov y V. Volkova, \emph{In hot pursuit of a stable wormhole in beyond Horndeski theory}, \href{https://arxiv.org/abs/2212.05969}{arXiv:2212.05969}.

\bibitem{MironovVolkova} S. Mironov y V. Volkova, \emph{Complete stability for spherically symmetric backgrounds in beyond Horndeski theory}, \href{https://arxiv.org/abs/2404.06297}{arXiv:2404.06297}.

\bibitem{repo} \emph{gusanos: híbrido estrella--agujero de gusano} (2026), \url{https://github.com/iivanlopez/gusanos}.

\bibitem{Kobayashi} T. Kobayashi, \emph{Horndeski theory and beyond: a review}, Rept. Prog. Phys. \textbf{82} (2019) 086901, \href{https://arxiv.org/abs/1901.07183}{arXiv:1901.07183}.

\bibitem{KMS1} T. Kobayashi, H. Motohashi y T. Suyama, \emph{Black hole perturbation in the most general scalar-tensor theory with second-order field equations I: The odd-parity sector}, Phys. Rev. D \textbf{85} (2012) 084025, \href{https://arxiv.org/abs/1202.4893}{arXiv:1202.4893}.

\bibitem{KMS2} T. Kobayashi, H. Motohashi y T. Suyama, \emph{Black hole perturbation in the most general scalar-tensor theory with second-order field equations II: the even-parity sector}, Phys. Rev. D \textbf{89} (2014) 084042, \href{https://arxiv.org/abs/1402.6740}{arXiv:1402.6740}.

\bibitem{Heisenberg} L. Heisenberg, \emph{Generalization of the Proca Action}, JCAP \textbf{05} (2014) 015, \href{https://arxiv.org/abs/1402.7026}{arXiv:1402.7026}.

\bibitem{LuPang} H. Lü y Y. Pang, \emph{Horndeski gravity as $D\to4$ limit of Gauss--Bonnet}, Phys. Lett. B \textbf{809} (2020) 135717, \href{https://arxiv.org/abs/2003.11552}{arXiv:2003.11552}.

\bibitem{BabichevCharmousis} E. Babichev y C. Charmousis, \emph{Dressing a black hole with a time-dependent Galileon}, JHEP \textbf{08} (2014) 106, \href{https://arxiv.org/abs/1312.3204}{arXiv:1312.3204}.

\bibitem{BCK} A. Bakopoulos, C. Charmousis y P. Kanti, \emph{Traversable wormholes in beyond Horndeski theories}, JCAP \textbf{05} (2022) 022, \href{https://arxiv.org/abs/2111.09857}{arXiv:2111.09857}.

\bibitem{BCEL1} E. Babichev, C. Charmousis, G. Esposito-Farèse y A. Lehébel, \emph{Stability of a black hole and the speed of gravity waves within self-tuning cosmological models}, Phys. Rev. Lett. \textbf{120} (2018) 241101, \href{https://arxiv.org/abs/1712.04398}{arXiv:1712.04398}.

\bibitem{BCEL2} E. Babichev, C. Charmousis, G. Esposito-Farèse y A. Lehébel, \emph{Hamiltonian vs stability and application to Horndeski theory}, Phys. Rev. D \textbf{98} (2018) 104050, \href{https://arxiv.org/abs/1803.11444}{arXiv:1803.11444}.

\bibitem{OgawaKobayashiSuyama} H. Ogawa, T. Kobayashi y T. Suyama, \emph{Instability of hairy black holes in shift-symmetric Horndeski theories}, Phys. Rev. D \textbf{93} (2016) 064078, \href{https://arxiv.org/abs/1510.07400}{arXiv:1510.07400}.

\bibitem{deRhamZhang} C. de Rham y J. Zhang, \emph{Perturbations of stealth black holes in DHOST theories}, Phys. Rev. D \textbf{100} (2019) 124023, \href{https://arxiv.org/abs/1907.00699}{arXiv:1907.00699}.

\bibitem{HKMT} L. Heisenberg, R. Kase, M. Minamitsuji y S. Tsujikawa, \emph{Hairy black-hole solutions in generalized Proca theories}, Phys. Rev. D \textbf{96} (2017) 084049, \href{https://arxiv.org/abs/1705.09662}{arXiv:1705.09662}.

\bibitem{KaseOdd} R. Kase, M. Minamitsuji, S. Tsujikawa y Y.-l. Zhang, \emph{Black hole perturbations in vector-tensor theories: The odd-mode analysis}, JCAP \textbf{02} (2018) 048, \href{https://arxiv.org/abs/1801.01787}{arXiv:1801.01787}.

\bibitem{DeFeliceProca} A. De Felice, L. Heisenberg, R. Kase, S. Mukohyama, S. Tsujikawa y Y.-l. Zhang, \emph{Cosmology in generalized Proca theories}, JCAP \textbf{06} (2016) 048, \href{https://arxiv.org/abs/1603.05806}{arXiv:1603.05806}.

\bibitem{DeFeliceCubic} A. De Felice, R. Kase y S. Tsujikawa, \emph{Scrutinizing black hole stability in cubic vector Galileon theories}, \href{https://arxiv.org/abs/2409.15606}{arXiv:2409.15606}.

\bibitem{GarciaSaenz} S. Garcia-Saenz, A. Held y J. Zhang, \emph{Destabilization of black holes and stars by generalized Proca fields}, Phys. Rev. Lett. \textbf{127} (2021) 131104, \href{https://arxiv.org/abs/2104.08049}{arXiv:2104.08049}.

\bibitem{Ganiyeva} N. Ganiyeva, B.~J. Barros, Á. de la Cruz-Dombriz y F.~S.~N. Lobo, \emph{Vector wormholes as conduits for matter interaction}, Eur. Phys. J. C \textbf{85} (2025) 914, \href{https://arxiv.org/abs/2412.19408}{arXiv:2412.19408}.

\bibitem{KarloviniSamuelsson} M. Karlovini y L. Samuelsson, \emph{Elastic stars in general relativity: I. Foundations and equilibrium models}, Class. Quant. Grav. \textbf{20} (2003) 3613, \href{https://arxiv.org/abs/gr-qc/0211026}{arXiv:gr-qc/0211026}.

\bibitem{KKK} P. Kanti, B. Kleihaus y J. Kunz, \emph{Wormholes in dilatonic Einstein--Gauss--Bonnet theory}, Phys. Rev. Lett. \textbf{107} (2011) 271101, \href{https://arxiv.org/abs/1108.3003}{arXiv:1108.3003}.

\bibitem{Canate} P. Cañate, J. Sultana y D. Kazanas, \emph{Ellis wormhole without a phantom scalar field}, Phys. Rev. D \textbf{100} (2019) 064007, \href{https://arxiv.org/abs/1907.09463}{arXiv:1907.09463}.

\end{thebibliography}
\end{document}
